% GENERATED FILE. Edit canonical lessons, then run npm run generate:books.
% canonical-source-hash: fnv1a64:3d826d4d13df0544
% canonical-lessons: ZH-C390-queding, ZH-C390-yueding, ZH-W390-ding-settle, ZH-C390-dingqi, ZH-C390-ruqi, ZH-C390-dayue

\chapter{As We Agreed}
\label{ch:zh-as-we-agreed}
\hlchaptermodality{390}

\begin{chapteropening}
\textbf{By the end of this chapter:} \emph{I can talk about being sure, agreeing on a time, doing things regularly, keeping to a date and giving a rough time.}

Complete the last lesson of chapter 390: I can talk about being sure, agreeing on a time, doing things regularly, keeping to a date and giving a rough time.
\end{chapteropening}

\section[quèdìng]{quèdìng — to be certain; to settle on}

\label{lesson:ZH-C390-queding}



\begin{quote}
Before the new one: say the Mandarin for as a precaution, then the Mandarin for to fall due.
\end{quote}

\subsection*{You'll want to know: quèdìng}
\emph{quèdìng} — \textquotedblleft{}to be certain; to settle on\textquotedblright{}.

\begin{quote}
\emph{wǒ hái bú quèdìng míngtiān néng bu néng qù} — I'm not sure yet whether I can go tomorrow.
\end{quote}

Its characters come later; for now it is a word for the ear.

\begin{cousinweb}
\emph{quèdìng} is built from \emph{què} 'firm' and \emph{dìng} 'settled', both inherited from Middle Chinese, \emph{khaewk} and \emph{dengH} (Guangyun, in Baxter's transcription), which glosses the first 'steadfast' and the second 'calm'.
\end{cousinweb}

\begin{grammarlens}[title={What you've built}]
One more word for saying what would happen.
\end{grammarlens}

\subsection*{Your turn}
\begin{itemize}
\raggedright
  \item \emph{Say it:} \emph{quèdìng}
  \item \emph{Say it:} \emph{quèdìng}, once more
  \item \emph{Say it:} \emph{wǒ hái bú quèdìng míngtiān néng bu néng qù} — I'm not sure yet whether I can go tomorrow.
  \item \emph{Recall:} say the Mandarin for as a precaution, then the Mandarin for to fall due, then say \emph{quèdìng} again
\end{itemize}

\begin{tcolorbox}[breakable,colback=teal!4,colframe=teal!35!black,title={Before you move on}]
What does quèdìng mean? (\textquotedblleft{}To be certain; to settle on\textquotedblright{}.) Say it once more.
\end{tcolorbox}
\section[\texorpdfstring{\zh{约定} (yuēdìng)}{约定 (yuēdìng)}]{\zh{约定} (yuēdìng) — to agree on a time; to fix (a time or place) together}

\label{lesson:ZH-C390-yueding}



\begin{quote}
Before the new one: say the Mandarin for to fall due, then the Mandarin for to be certain.
\end{quote}

\subsection*{You'll want to know: \zh{约定}}
\textbf{\zh{约定}} — \emph{yuēdìng} — \textquotedblleft{}to agree on a time; to fix (a time or place) together\textquotedblright{}.

\begin{quote}
\textbf{\zh{我们约定明天早上八点见面}} \emph{wǒmen yuēdìng míngtiān zǎoshang bā diǎn jiànmiàn} — We agreed to meet at eight tomorrow morning.
\end{quote}

The new character, in pieces: \zh{定} — eight strokes in two pieces: the roof of \zh{安} on top, and under it the five strokes at the foot of \zh{是}. The next lesson puts it in your hand.

\begin{cousinweb}
\textbf{\zh{约定}} (\emph{yuēdìng}) is built from \zh{约} 'to bind' and \zh{定} 'settled', both inherited from Middle Chinese, \emph{'jak} and \emph{dengH} (Guangyun, in Baxter's transcription), which glosses the second 'calm'.
\end{cousinweb}

\begin{grammarlens}[title={What you've built}]
One more word for saying what would happen.
\end{grammarlens}

\subsection*{Your turn}
\begin{itemize}
\raggedright
  \item \emph{Say it:} \emph{yuēdìng}
  \item \emph{Say it:} \emph{yuēdìng}, once more
  \item \emph{Say it:} \emph{wǒmen yuēdìng míngtiān zǎoshang bā diǎn jiànmiàn} — We agreed to meet at eight tomorrow morning.
  \item \emph{Recall:} say the Mandarin for to fall due, then the Mandarin for to be certain, then say \emph{yuēdìng} again
\end{itemize}

\begin{tcolorbox}[breakable,colback=teal!4,colframe=teal!35!black,title={Before you move on}]
What does \zh{约定} mean? (\textquotedblleft{}To agree on a time; to fix (a time or place) together\textquotedblright{}.) Say it once more.
\end{tcolorbox}
\section[\texorpdfstring{\zh{定} (dìng)}{定 (dìng)}]{\zh{定} — the roof of \zh{安} over the foot of \zh{是}}

\label{lesson:ZH-W390-ding-settle}



\begin{quote}
The word first, then the character inside it.

\begin{itemize}
\raggedright
  \item \emph{Recall:} say \emph{yuēdìng} — \textbf{R1}, one lesson back, before it decays
\end{itemize}

\begin{quote}
\zh{安} \zh{是}
\end{quote}

\emph{Read it:} each one aloud, with nothing said first — \emph{ān}, \emph{shì}
\end{quote}

\subsection*{Script — the roof of \zh{安} over the foot of \zh{是}}
\hlblockfigure{\detokenize{figures/ZH-W390-ding-settle-filmstrip.pdf}}{How \zh{定} is written, stroke by stroke}

\begin{quote}
\zh{定}
\end{quote}

\textbf{Two pieces, top before bottom.}

\textbf{The roof, three strokes,} as on top of \zh{安}: a dot, a short stroke down on the left, and a bar that turns down and hooks.

\textbf{Under it, five strokes, the foot of \zh{是}.} A bar. A short vertical down from its middle. A short bar to the right of the vertical. Then a stroke falling to the left, and a long stroke falling to the right that sweeps along the foot.

3 + 5 = eight.

\subsection*{Your turn}
\begin{itemize}
\raggedright
  \item \emph{Write it:} \zh{定}, all eight, with the model in front of you
  \item \emph{Cover:} write it again from memory, and count the pieces, not the strokes
\end{itemize}

\begin{tcolorbox}[breakable,colback=teal!4,colframe=teal!35!black,title={Before you move on}]
How many strokes in \zh{定}? (\textbf{Eight}.) What is on top? (\textbf{The roof of \zh{安}}.) Where have you made the bottom piece before? (\textbf{At the foot of \zh{是}}.)

Source: \href{https://raw.githubusercontent.com/chanind/hanzi-writer-data/68d10a4b21150cae5e1ebbd223eed289cf32d90c/data/\%E5\%AE\%9A.json}{Hanzi Writer Data \zh{定}, pinned stroke paths and medians}.
\end{tcolorbox}
\section[\texorpdfstring{\zh{定期} (dìngqī)}{定期 (dìngqī)}]{\zh{定期} (dìngqī) — at regular intervals; periodically}

\label{lesson:ZH-C390-dingqi}



\begin{quote}
Before the new one: say the Mandarin for to be certain, then the Mandarin for to agree on a time.
\end{quote}

\subsection*{You'll want to know: \zh{定期}}
\textbf{\zh{定期}} — \emph{dìngqī} — \textquotedblleft{}at regular intervals; periodically\textquotedblright{}.

\begin{quote}
\textbf{\zh{你最好定期去看医生}} \emph{nǐ zuìhǎo dìngqī qù kàn yīshēng} — You'd better see a doctor regularly.
\end{quote}

\begin{cousinweb}
\textbf{\zh{定期}} (\emph{dìngqī}) is built from \zh{定} 'settled' and \zh{期} 'a set time', both inherited from Middle Chinese, \emph{dengH} and \emph{gi} (Guangyun, in Baxter's transcription), which glosses the second 'a limit'. The level tone of qī is irregular, as in \zh{星期}.
\end{cousinweb}

\begin{grammarlens}[title={What you've built}]
One more word for saying what would happen.
\end{grammarlens}

\subsection*{Your turn}
\begin{itemize}
\raggedright
  \item \emph{Say it:} \emph{dìngqī}
  \item \emph{Say it:} \emph{dìngqī}, once more
  \item \emph{Say it:} \emph{nǐ zuìhǎo dìngqī qù kàn yīshēng} — You'd better see a doctor regularly.
  \item \emph{Recall:} say the Mandarin for to be certain, then the Mandarin for to agree on a time, then say \emph{dìngqī} again
\end{itemize}

\begin{tcolorbox}[breakable,colback=teal!4,colframe=teal!35!black,title={Before you move on}]
What does \zh{定期} mean? (\textquotedblleft{}At regular intervals; periodically\textquotedblright{}.) Say it once more.
\end{tcolorbox}
\section[\texorpdfstring{\zh{如期} (rúqī)}{如期 (rúqī)}]{\zh{如期} (rúqī) — as planned (on the agreed date)}

\label{lesson:ZH-C390-ruqi}



\begin{quote}
Before the new one: say the Mandarin for to agree on a time, then the Mandarin for at regular intervals.
\end{quote}

\subsection*{You'll want to know: \zh{如期}}
\textbf{\zh{如期}} — \emph{rúqī} — \textquotedblleft{}as planned (on the agreed date)\textquotedblright{}.

\begin{quote}
\textbf{\zh{要是不下雨我们就如期去山里}} \emph{yàoshi bú xiàyǔ wǒmen jiù rúqī qù shān lǐ} — If it doesn't rain, we'll go to the hills as planned.
\end{quote}

\begin{cousinweb}
\textbf{\zh{如期}} (\emph{rúqī}) is built from \zh{如} 'as' and \zh{期} 'a set time', both inherited from Middle Chinese, \emph{nyo} and \emph{gi} (Guangyun, in Baxter's transcription), which glosses the first 'like' and the second 'an appointed meeting'.
\end{cousinweb}

\begin{grammarlens}[title={What you've built}]
One more word for saying what would happen.
\end{grammarlens}

\subsection*{Your turn}
\begin{itemize}
\raggedright
  \item \emph{Say it:} \emph{rúqī}
  \item \emph{Say it:} \emph{rúqī}, once more
  \item \emph{Say it:} \emph{yàoshi bú xiàyǔ wǒmen jiù rúqī qù shān lǐ} — If it doesn't rain, we'll go to the hills as planned.
  \item \emph{Recall:} say the Mandarin for to agree on a time, then the Mandarin for at regular intervals, then say \emph{rúqī} again
\end{itemize}

\begin{tcolorbox}[breakable,colback=teal!4,colframe=teal!35!black,title={Before you move on}]
What does \zh{如期} mean? (\textquotedblleft{}As planned (on the agreed date)\textquotedblright{}.) Say it once more.
\end{tcolorbox}
\section[\texorpdfstring{\zh{大约} (dàyuē)}{大约 (dàyuē)}]{\zh{大约} (dàyuē) — approximately; roughly (a time or number)}

\label{lesson:ZH-C390-dayue}



\begin{quote}
Before the new one: say the Mandarin for at regular intervals, then the Mandarin for on the agreed date.
\end{quote}

\subsection*{You'll want to know: \zh{大约}}
\textbf{\zh{大约}} — \emph{dàyuē} — \textquotedblleft{}approximately; roughly (a time or number)\textquotedblright{}.

\begin{quote}
\emph{yàoshi bù dǔchē wǒ dàyuē bā diǎn dào} — If there's no traffic, I'll be there at about eight.
\end{quote}

\begin{cousinweb}
\textbf{\zh{大约}} (\emph{dàyuē}) is built from \zh{大} 'big' and \zh{约} 'to bind', both inherited from Middle Chinese, \emph{dajH} and \emph{'jak} (Guangyun, in Baxter's transcription), which glosses the second 'to tie up' and 'sparing'. How the pair came to mean 'roughly' is not settled.
\end{cousinweb}

\begin{grammarlens}[title={What you've built}]
One more word for saying what would happen.
\end{grammarlens}

\subsection*{Your turn}
\begin{itemize}
\raggedright
  \item \emph{Say it:} \emph{dàyuē}
  \item \emph{Say it:} \emph{dàyuē}, once more
  \item \emph{Say it:} \emph{yàoshi bù dǔchē wǒ dàyuē bā diǎn dào} — If there's no traffic, I'll be there at about eight.
  \item \emph{Recall:} say the Mandarin for at regular intervals, then the Mandarin for on the agreed date, then say \emph{dàyuē} again
\end{itemize}

\begin{tcolorbox}[breakable,colback=teal!4,colframe=teal!35!black,title={Before you move on}]
What does \zh{大约} mean? (\textquotedblleft{}Approximately; roughly (a time or number)\textquotedblright{}.) Say it once more.
\end{tcolorbox}
